%% Journal information
%% Supplied to authors by publisher for camera-ready submission;
%% use defaults for review submission.
\acmConference[IFL'19]{International Symposium on Implementation and Application of Functional Languages}{September 2019}{Singapore}
\acmYear{2020}
\copyrightyear{2020}
\acmDOI{} % \acmDOI{10.1145/nnnnnnn.nnnnnnn}
\startPage{1}

%% Copyright information
%% Supplied to authors (based on authors' rights management selection;
%% see authors.acm.org) by publisher for camera-ready submission;
%% use 'none' for review submission.
\setcopyright{none}
%\setcopyright{acmcopyright}
%\setcopyright{acmlicensed}
%\setcopyright{rightsretained}
%\copyrightyear{2018}           %% If different from \acmYear

%% Bibliography style
\bibliographystyle{ACM-Reference-Format}
%% Citation style
%% Note: author/year citations are required for papers published as an
%% issue of PACMPL.
\citestyle{acmauthoryear}   %% For author/year citations


%%%%%%%%%%%%%%%%%%%%%%%%%%%%%%%%%%%%%%%%%%%%%%%%%%%%%%%%%%%%%%%%%%%%%%
%% Note: Authors migrating a paper from PACMPL format to traditional
%% SIGPLAN proceedings format must update the '\documentclass' and
%% topmatter commands above; see 'acmart-sigplanproc-template.tex'.
%%%%%%%%%%%%%%%%%%%%%%%%%%%%%%%%%%%%%%%%%%%%%%%%%%%%%%%%%%%%%%%%%%%%%%


\usepackage{mathpartir} % \infer
\usepackage{listings}

\lstdefinelanguage{IR} {
mathescape=true,
texcl=false,
morekeywords=[1] {
let, ret, proj, case, of, pap, ctor, inc, dec, reset, reuse, in
},
literate=
{→}{{\ensuremath{\rightarrow\,}}}1
{->}{{\ensuremath{\rightarrow\,}}}1
{_i}{{\ensuremath{_i}}}1
{_1}{{\ensuremath{_1}}}1
{_2}{{\ensuremath{_2}}}1
{_3}{{\ensuremath{_3}}}1
{_4}{{\ensuremath{_4}}}1
{_5}{{\ensuremath{_5}}}1
{₁}{{\ensuremath{_1}}}1
{₂}{{\ensuremath{_2}}}1
{₃}{{\ensuremath{_3}}}1
{₄}{{\ensuremath{_4}}}1
{₅}{{\ensuremath{_5}}}1
{:}{{\ensuremath{\,}}:{\ensuremath{\,}}}1
{;}{{\ensuremath{\,}};}1
{collecto}{{\ensuremath{\textit{collect}_\own{}}}}1
{₆}{{\ensuremath{_6}}}1,
%%{collect_o}{{\ensuremath{\textit{collect}_\own{}}}}1,
%% basicstyle=\small,
identifierstyle={\rmfamily\itshape\color{black}},
keywordstyle=[1]{\ttfamily\bfseries\color{black}},
stringstyle=\ttfamily,
columns=fullflexible
}

\lstset{language=IR}
\usepackage{amsthm} % \newtheorem
\usepackage{amssymb} % \twoheadrightarrow
\usepackage[capitalise]{cleveref} % \cref
\usepackage{hyperref} % \url

\newtheorem{lem}{Lemma}
\newtheorem{thm}{Theorem}
\newtheorem{cor}{Corollary}
\newtheorem{defn}{Definition}
\DeclareMathOperator{\dom}{dom}
\DeclareMathOperator{\ran}{ran}
\DeclareMathOperator{\dec}{dec}
\DeclareMathOperator{\inc}{inc}
\DeclareMathOperator{\live}{live}
\DeclareMathOperator{\vars}{vars}
\DeclareMathOperator{\erase}{erase}
\DeclareMathOperator{\FV}{FV}
\DeclareMathOperator{\valof}{valof}
\DeclareMathOperator{\reachable}{reachable}
\newcommand\RC{_{\textit{RC}}}
\newcommand\rcIR{\ensuremath{\lambda\RC}}
\def\Reset#1{\kw{reset}\ #1}
\def\Reuse_#1#2#3{\kw{reuse}\ #2\ \kw{in}\ \kw{ctor}_#1\ #3}
\def\Inc#1;#2{\kw{inc}\ #1;\ #2}
\def\Dec#1;#2{\kw{dec}\ #1;\ #2}

\newcommand{\var}[1]{\texttt{#1}}
\newcommand\abs[1]{\mid #1 \mid}

\newcommand\pure{_{\textit{pure}}}
\newcommand\pureIR{\ensuremath{\lambda\pure}}
\newcommand\psem[3]{#1 \vdash #2 \Downarrow #3}
\newcommand\ty[1]{\mathit{#1}}
\newcommand\syncat[1]{&\in \ty{#1}}
\newcommand\kw[1]{\texttt{\textbf{#1}}}
\newcommand\many[1]{\overline{#1}}
\newcommand\rcsem[5]{#1 \vdash \langle #2, #3 \rangle \Downarrow \langle #4, #5 \rangle}
\newcommand\nullptr{\rule{0.5em}{0.5em}}

\def\Pap#1#2{\kw{pap}\ #1\ #2}
\def\Ctor_#1#2{\kw{ctor}_#1\ #2}
\def\Proj_#1#2{\kw{proj}_#1\ #2}
\def\Ret#1{\kw{ret}\ #1}
\def\Let#1=#2;#3{\kw{let}\ #1 = #2;\ #3}
\def\Case#1#2{\kw{case}\ #1\ \kw{of}\ #2}
\def\Lambda#1.#2{\lambda\ #1.\ #2}

\newcommand\own{\ensuremath{\mathbb{O}}}
\newcommand\bor{\ensuremath{\mathbb{B}}}
\newcommand\reset{\ensuremath{\mathbb{R}}}
\newcommand\reuse{_{\textit{reuse}}}
\newcommand\squote{\textnormal{\textquotesingle}}

%% Title information
\title{Counting Immutable Beans}
\subtitle{Reference Counting Optimized for Purely Functional Programming}
%\titlenote{with title note}             %% \titlenote is optional;
%                                        %% can be repeated if necessary;
%                                        %% contents suppressed with 'anonymous'
%\subtitle{Subtitle}                     %% \subtitle is optional
%\subtitlenote{with subtitle note}       %% \subtitlenote is optional;
%                                        %% can be repeated if necessary;
%                                        %% contents suppressed with 'anonymous'


%% Author information
%% Contents and number of authors suppressed with 'anonymous'.
%% Each author should be introduced by \author, followed by
%% \authornote (optional), \orcid (optional), \affiliation, and
%% \email.
%% An author may have multiple affiliations and/or emails; repeat the
%% appropriate command.
%% Many elements are not rendered, but should be provided for metadata
%% extraction tools.

%% Author with single affiliation.
\author{Sebastian Ullrich}
\affiliation{
  %\position{Position1}
  %\department{Department of Informatics}
  \institution{Karlsruhe Institute of Technology}
  %\streetaddress{Street1 Address1}
  %\city{City1}
  %\state{State1}
  %\postcode{Post-Code1}
  \country{Germany}
}
\email{sebastian.ullrich@kit.edu}

%% Author with two affiliations and emails.
\author{Leonardo de Moura}
\affiliation{
  % \position{Position2a}
  % \department{Department2a}             %% \department is recommended
  \institution{Microsoft Research}           %% \institution is required
  % \streetaddress{Street2a Address2a}
  % \city{City2a}
  %\state{State2a}
  %\postcode{Post-Code2a}
  \country{USA}                   %% \country is recommended
}
\email{leonardo@microsoft.com}

%% Abstract
%% Note: \begin{abstract}...\end{abstract} environment must come
%% before \maketitle command
\begin{abstract}
  Most functional languages rely on some kind of garbage collection
  for automatic memory management. They usually eschew reference
  counting in favor of a tracing garbage collector, which has less
  bookkeeping overhead at runtime. On the other hand, having
  an exact reference count of each value can enable
  optimizations such as destructive updates. We explore these
  optimization opportunities in the context of an eager, purely
  functional programming language. We propose a new mechanism for efficiently reclaiming
  memory used by nonshared values, reducing stress on the
  global memory allocator. We describe an approach for
  minimizing the number of reference counts updates using borrowed
  references and a heuristic for automatically inferring borrow
  annotations. We implemented all these techniques in a
  new compiler for an eager and purely functional programming language
  with support for multi-threading. Our preliminary experimental
  results demonstrate our approach is competitive and often
  outperforms state-of-the-art compilers.
\end{abstract}


%% 2012 ACM Computing Classification System (CSS) concepts
%% Generate at 'http://dl.acm.org/ccs/ccs.cfm'.
\begin{CCSXML}
<ccs2012>
<concept>
<concept_id>10011007.10011006.10011041.10011048</concept_id>
<concept_desc>Software and its engineering~Runtime environments</concept_desc>
<concept_significance>500</concept_significance>
</concept>
<concept>
<concept_id>10011007.10010940.10010941.10010949.10010950.10010954</concept_id>
<concept_desc>Software and its engineering~Garbage collection</concept_desc>
<concept_significance>500</concept_significance>
</concept>
</ccs2012>
\end{CCSXML}

\ccsdesc[500]{Software and its engineering~Runtime environments}
\ccsdesc[500]{Software and its engineering~Garbage collection}
%% End of generated code


%% Keywords
%% comma separated list
\keywords{purely functional programming, reference counting, Lean}  %% \keywords are mandatory in final camera-ready submission
